% ----------------------------------------------------
%
%          Desugar to ICFP lambdas
%
%         Was: desugaring.ltx
%
% ----------------------------------------------------

%% ODER: format ==         = "\mathrel{==}"
%% ODER: format /=         = "\neq "
%
%
\makeatletter
\@ifundefined{lhs2tex.lhs2tex.sty.read}%
  {\@namedef{lhs2tex.lhs2tex.sty.read}{}%
   \newcommand\SkipToFmtEnd{}%
   \newcommand\EndFmtInput{}%
   \long\def\SkipToFmtEnd#1\EndFmtInput{}%
  }\SkipToFmtEnd

\newcommand\ReadOnlyOnce[1]{\@ifundefined{#1}{\@namedef{#1}{}}\SkipToFmtEnd}
\usepackage{amstext}
\usepackage{amssymb}
\usepackage{stmaryrd}
\DeclareFontFamily{OT1}{cmtex}{}
\DeclareFontShape{OT1}{cmtex}{m}{n}
  {<5><6><7><8>cmtex8
   <9>cmtex9
   <10><10.95><12><14.4><17.28><20.74><24.88>cmtex10}{}
\DeclareFontShape{OT1}{cmtex}{m}{it}
  {<-> ssub * cmtt/m/it}{}
\newcommand{\texfamily}{\fontfamily{cmtex}\selectfont}
\DeclareFontShape{OT1}{cmtt}{bx}{n}
  {<5><6><7><8>cmtt8
   <9>cmbtt9
   <10><10.95><12><14.4><17.28><20.74><24.88>cmbtt10}{}
\DeclareFontShape{OT1}{cmtex}{bx}{n}
  {<-> ssub * cmtt/bx/n}{}
\newcommand{\tex}[1]{\text{\texfamily#1}}	% NEU

\newcommand{\Sp}{\hskip.33334em\relax}


\newcommand{\Conid}[1]{\mathit{#1}}
\newcommand{\Varid}[1]{\mathit{#1}}
\newcommand{\anonymous}{\kern0.06em \vbox{\hrule\@width.5em}}
\newcommand{\plus}{\mathbin{+\!\!\!+}}
\newcommand{\bind}{\mathbin{>\!\!\!>\mkern-6.7mu=}}
\newcommand{\rbind}{\mathbin{=\mkern-6.7mu<\!\!\!<}}% suggested by Neil Mitchell
\newcommand{\sequ}{\mathbin{>\!\!\!>}}
\renewcommand{\leq}{\leqslant}
\renewcommand{\geq}{\geqslant}
\usepackage{polytable}

%mathindent has to be defined
\@ifundefined{mathindent}%
  {\newdimen\mathindent\mathindent\leftmargini}%
  {}%

\def\resethooks{%
  \global\let\SaveRestoreHook\empty
  \global\let\ColumnHook\empty}
\newcommand*{\savecolumns}[1][default]%
  {\g@addto@macro\SaveRestoreHook{\savecolumns[#1]}}
\newcommand*{\restorecolumns}[1][default]%
  {\g@addto@macro\SaveRestoreHook{\restorecolumns[#1]}}
\newcommand*{\aligncolumn}[2]%
  {\g@addto@macro\ColumnHook{\column{#1}{#2}}}

\resethooks

\newcommand{\onelinecommentchars}{\quad-{}- }
\newcommand{\commentbeginchars}{\enskip\{-}
\newcommand{\commentendchars}{-\}\enskip}

\newcommand{\visiblecomments}{%
  \let\onelinecomment=\onelinecommentchars
  \let\commentbegin=\commentbeginchars
  \let\commentend=\commentendchars}

\newcommand{\invisiblecomments}{%
  \let\onelinecomment=\empty
  \let\commentbegin=\empty
  \let\commentend=\empty}

\visiblecomments

\newlength{\blanklineskip}
\setlength{\blanklineskip}{0.66084ex}

\newcommand{\hsindent}[1]{\quad}% default is fixed indentation
\let\hspre\empty
\let\hspost\empty
\newcommand{\NB}{\textbf{NB}}
\newcommand{\Todo}[1]{$\langle$\textbf{To do:}~#1$\rangle$}

\EndFmtInput
\makeatother
%
%
%
%
%
%
% This package provides two environments suitable to take the place
% of hscode, called "plainhscode" and "arrayhscode". 
%
% The plain environment surrounds each code block by vertical space,
% and it uses \abovedisplayskip and \belowdisplayskip to get spacing
% similar to formulas. Note that if these dimensions are changed,
% the spacing around displayed math formulas changes as well.
% All code is indented using \leftskip.
%
% Changed 19.08.2004 to reflect changes in colorcode. Should work with
% CodeGroup.sty.
%
\ReadOnlyOnce{polycode.fmt}%
\makeatletter

\newcommand{\hsnewpar}[1]%
  {{\parskip=0pt\parindent=0pt\par\vskip #1\noindent}}

% can be used, for instance, to redefine the code size, by setting the
% command to \small or something alike
\newcommand{\hscodestyle}{}

% The command \sethscode can be used to switch the code formatting
% behaviour by mapping the hscode environment in the subst directive
% to a new LaTeX environment.

\newcommand{\sethscode}[1]%
  {\expandafter\let\expandafter\hscode\csname #1\endcsname
   \expandafter\let\expandafter\endhscode\csname end#1\endcsname}

% "compatibility" mode restores the non-polycode.fmt layout.

\newenvironment{compathscode}%
  {\par\noindent
   \advance\leftskip\mathindent
   \hscodestyle
   \let\\=\@normalcr
   \let\hspre\(\let\hspost\)%
   \pboxed}%
  {\endpboxed\)%
   \par\noindent
   \ignorespacesafterend}

\newcommand{\compaths}{\sethscode{compathscode}}

% "plain" mode is the proposed default.
% It should now work with \centering.
% This required some changes. The old version
% is still available for reference as oldplainhscode.

\newenvironment{plainhscode}%
  {\hsnewpar\abovedisplayskip
   \advance\leftskip\mathindent
   \hscodestyle
   \let\hspre\(\let\hspost\)%
   \pboxed}%
  {\endpboxed%
   \hsnewpar\belowdisplayskip
   \ignorespacesafterend}

\newenvironment{oldplainhscode}%
  {\hsnewpar\abovedisplayskip
   \advance\leftskip\mathindent
   \hscodestyle
   \let\\=\@normalcr
   \(\pboxed}%
  {\endpboxed\)%
   \hsnewpar\belowdisplayskip
   \ignorespacesafterend}

% Here, we make plainhscode the default environment.

\newcommand{\plainhs}{\sethscode{plainhscode}}
\newcommand{\oldplainhs}{\sethscode{oldplainhscode}}
\plainhs

% The arrayhscode is like plain, but makes use of polytable's
% parray environment which disallows page breaks in code blocks.

\newenvironment{arrayhscode}%
  {\hsnewpar\abovedisplayskip
   \advance\leftskip\mathindent
   \hscodestyle
   \let\\=\@normalcr
   \(\parray}%
  {\endparray\)%
   \hsnewpar\belowdisplayskip
   \ignorespacesafterend}

\newcommand{\arrayhs}{\sethscode{arrayhscode}}

% The mathhscode environment also makes use of polytable's parray 
% environment. It is supposed to be used only inside math mode 
% (I used it to typeset the type rules in my thesis).

\newenvironment{mathhscode}%
  {\parray}{\endparray}

\newcommand{\mathhs}{\sethscode{mathhscode}}

% texths is similar to mathhs, but works in text mode.

\newenvironment{texthscode}%
  {\(\parray}{\endparray\)}

\newcommand{\texths}{\sethscode{texthscode}}

% The framed environment places code in a framed box.

\def\codeframewidth{\arrayrulewidth}
\RequirePackage{calc}

\newenvironment{framedhscode}%
  {\parskip=\abovedisplayskip\par\noindent
   \hscodestyle
   \arrayrulewidth=\codeframewidth
   \tabular{@{}|p{\linewidth-2\arraycolsep-2\arrayrulewidth-2pt}|@{}}%
   \hline\framedhslinecorrect\\{-1.5ex}%
   \let\endoflinesave=\\
   \let\\=\@normalcr
   \(\pboxed}%
  {\endpboxed\)%
   \framedhslinecorrect\endoflinesave{.5ex}\hline
   \endtabular
   \parskip=\belowdisplayskip\par\noindent
   \ignorespacesafterend}

\newcommand{\framedhslinecorrect}[2]%
  {#1[#2]}

\newcommand{\framedhs}{\sethscode{framedhscode}}

% The inlinehscode environment is an experimental environment
% that can be used to typeset displayed code inline.

\newenvironment{inlinehscode}%
  {\(\def\column##1##2{}%
   \let\>\undefined\let\<\undefined\let\\\undefined
   \newcommand\>[1][]{}\newcommand\<[1][]{}\newcommand\\[1][]{}%
   \def\fromto##1##2##3{##3}%
   \def\nextline{}}{\) }%

\newcommand{\inlinehs}{\sethscode{inlinehscode}}

% The joincode environment is a separate environment that
% can be used to surround and thereby connect multiple code
% blocks.

\newenvironment{joincode}%
  {\let\orighscode=\hscode
   \let\origendhscode=\endhscode
   \def\endhscode{\def\hscode{\endgroup\def\@currenvir{hscode}\\}\begingroup}
   %\let\SaveRestoreHook=\empty
   %\let\ColumnHook=\empty
   %\let\resethooks=\empty
   \orighscode\def\hscode{\endgroup\def\@currenvir{hscode}}}%
  {\origendhscode
   \global\let\hscode=\orighscode
   \global\let\endhscode=\origendhscode}%

\makeatother
\EndFmtInput
%
%
%
% First, let's redefine the forall, and the dot.
%
%
% This is made in such a way that after a forall, the next
% dot will be printed as a period, otherwise the formatting
% of `comp_` is used. By redefining `comp_`, as suitable
% composition operator can be chosen. Similarly, period_
% is used for the period.
%
\ReadOnlyOnce{forall.fmt}%
\makeatletter

% The HaskellResetHook is a list to which things can
% be added that reset the Haskell state to the beginning.
% This is to recover from states where the hacked intelligence
% is not sufficient.

\let\HaskellResetHook\empty
\newcommand*{\AtHaskellReset}[1]{%
  \g@addto@macro\HaskellResetHook{#1}}
\newcommand*{\HaskellReset}{\HaskellResetHook}

\global\let\hsforallread\empty

\newcommand\hsforall{\global\let\hsdot=\hsperiodonce}
\newcommand*\hsperiodonce[2]{#2\global\let\hsdot=\hscompose}
\newcommand*\hscompose[2]{#1}

\AtHaskellReset{\global\let\hsdot=\hscompose}

% In the beginning, we should reset Haskell once.
\HaskellReset

\makeatother
\EndFmtInput
% ----------------------
%% Stuff for TimCore


%% 'g' inferences





%% Intersection of effects

%% --------- Effects ------------


%% ----------------------


%% --------------------------
%%        Desugaring rules
%% --------------------------



% ----------------------------------------------------
% New desugaring, just does the wrapping part:
% ----------------------------------------------------
%   Input for wrapping


%%format Hide (e) = "{\mathcal H}\llbracket{}" e "\rrbracket{}"

%% 'si' is short for "solve or infer" or "unification variable or not" resp

%% 'md' is short for "mode" or "unification variable or not" resp

%% flipsi is a no-op for now


%% EXX is just EX, but with no argument




%% ----------- Applications --------------
%% %format applyx a b = a "@" b
%% %format applyx a b = a "\," b

%% Function application with no arguments

%% Function application with 1 argument

%% Function application with 1 argument, no parens

%% Function application with 1 argument, without parenthesis

%% Function application with many arguments

%% ----------- End of Applications --------------





%% Sequences e1; ...; en; v

%% xs etc stole syntax used in examples.

%% Old version: | for BNF, [] for choice
%%%format choice = "[ \! ]"
%%%format `choice` = "\mathop{[ \! ]}"
%% 1mm is about 0.3em with the current font

%% New version: tall | for BNF, fat | for choice
%%format vbar = "\mathop{\bigm|}"
%%format choice = "\boldsymbol{|}"
%%format `choice` = "\mathop{\boldsymbol{|}}"




%% %format ok = "\textbf{ok}"

%% %format `eqsemi` = "\hspace{-0.5mm}\fatsemi{}"
%% %format eqsemiKW = "\fatsemi"

%% Arrays and tuples
%% %format arrKW = "\textbf{arr}"
%% %format arr (x) = arrKW "\{" x "\}"

%% %format vmap (x) = "\llangle" x "\rrangle"



%%  OLD SYNTAX   %format split (x) (y) (z) = "\textbf{split}(" x ")\{" y "," z "\}"
%% %format defKW = "\textbf{def}"
% only used in comments, but still prettier this way:
%%%%format map = "\textbf{map}"

%% -------Functions ---------------


%% functionOC has open/closed subscript


%% -------End of f unctions ---------------





%% Unification
%%%format == = "\!=\!"


% format := = "\mapsto"
% %format squiggt = "\leadsto"
% %format squiggt = "\!-\!\!>\!\!\!"



































%% We combined S and C into one


%% ----------- Metatheory --------------


%% ----------------------------
%%         Verification
%% ----------------------------

%% assert = succeed, abbreviated to suc



%% Effects checking, like succeeds{e}


%% <fx> brackets

\begin{figure}
$$
\begin{array}{@{}l}
\textbf{Syntax of Core Verse} \\
\begin{array}{lr@{\;\;}c@{\;\;}l}
 \text{Expressions} & e & ::= & \ensuremath{\Varid{x}} \vb \ensuremath{\Varid{k}} \vb \ensuremath{\Varid{o}} \vb \ensuremath{\Varid{e}_{\mathrm{1}} [\!\Varid{e}_{\mathrm{2}}]} \vb \ensuremath{\Varid{e}_{\mathrm{1}};\,\Varid{e}_{\mathrm{2}}}
     \vb \ensuremath{\Varid{e}_{\mathrm{1}}\hspace{-0.14em}=\hspace{-0.14em}\Varid{e}_{\mathrm{2}}} \vb \ensuremath{\exists{}\;\Varid{x}} \vb \ensuremath{\Varid{e}_{\mathrm{1}}\hspace{0.19em}\mathop{\rule[-0.1em]{0.15em}{0.75em}}\hspace{0.2em}\Varid{e}_{\mathrm{2}}}  \\
  && \vb &  \ensuremath{\langle\Varid{e}_{\mathrm{1}},..,\Varid{e}_{\Varid{n}}\rangle} \vb \ensuremath{\textbf{truth}\{\Varid{e}\}} \vb \ensuremath{\lambda\Varid{x} .\,\Varid{e}} \vb \ensuremath{\textbf{iter}(\Varid{v})\{\Varid{e}_{\mathrm{1}}\}\{\Varid{e}_{\mathrm{2}}\}}\\
  && \vb & \ensuremath{{\bf some}(\Varid{e})} \vb \ensuremath{\Varid{v}\;\!\blacktriangleright{}\!\;\Varid{e}} \vb \ensuremath{{\bf chk}\langle\omega\rangle\{\Varid{e}\}} \vb \ensuremath{\textbf{verify}(R;\,A)\{\Varid{e}\}} \\
\multicolumn{4}{l}{\hspace{3mm}\text{Note: ordinary \ensuremath{\lambda\Varid{x} .\,\Varid{e}}; no \ensuremath{\Varid{e}_{\mathrm{1}}\vartriangleright\Varid{e}_{\mathrm{2}}}; add \ensuremath{\textbf{verify}} and \ensuremath{{\bf some}} }} \\[2mm]
\end{array}\\

\mbox{{\bf Desugaring Mini Verse into Core Verse:} $\ensuremath{{\mathcal V}_{\sigma}\llbracket\Varid{t}\rrbracket{}\pi} = \ensuremath{\Varid{e}}$} \\

% -------------- Syntax ----------------
\multicolumn{1}{c}{\begin{array}{rcll}
  \ensuremath{\sigma} & ::= & \ensuremath{\textbf{x}} \vb \ensuremath{+} \vb \ensuremath{-}  & \text{``Mode'': execution, verify, check} \\
  \ensuremath{\pi} & ::= & \text{Set of in-scope lambda-bound variables} \\[1mm]
\end{array}} \\


% -------------- Basic rules ----------------
\begin{array}{lr@{\;\;}c@{\;\;}ll}
\rulename{mfail}  & \ensuremath{{\mathcal V}_{\sigma}\llbracket\textbf{fail}\rrbracket{}\pi}                  & = & \ensuremath{\textbf{fail}} \\
\rulename{mconst}  & \ensuremath{{\mathcal V}_{\sigma}\llbracket\Varid{k}\rrbracket{}\pi}                  & = & \ensuremath{\Varid{k}} \\
\rulename{mvar}    & \ensuremath{{\mathcal V}_{\sigma}\llbracket\Varid{x}\rrbracket{}\pi}                  & = & \ensuremath{\Varid{x}} \\
\rulename{mop}    & \ensuremath{{\mathcal V}_{\sigma}\llbracket\Varid{o}\rrbracket{}\pi}                  & = & \ensuremath{\Varid{o}} \\
\rulename{mapp}    & \ensuremath{{\mathcal V}_{\sigma}\llbracket\Varid{e}_{\mathrm{1}} [\!\Varid{e}_{\mathrm{2}}]\rrbracket{}\pi}
   & = & \ensuremath{({\mathcal V}_{\sigma}\llbracket\Varid{e}_{\mathrm{1}}\rrbracket{}\pi) [\!{\mathcal V}_{\sigma}\llbracket\Varid{e}_{\mathrm{2}}\rrbracket{}\pi]} \\
\rulename{mtup} & \ensuremath{{\mathcal V}_{\sigma}\llbracket\langle\Varid{e}_{\mathrm{1}},\mydots,\Varid{e}_{\Varid{n}}\rangle\rrbracket{}\pi}
    & = & \ensuremath{\langle{\mathcal V}_{\sigma}\llbracket\Varid{e}_{\mathrm{1}}\rrbracket{}\pi,\mydots,{\mathcal V}_{\sigma}\llbracket\Varid{e}_{\Varid{n}}\rrbracket{}\pi\rangle} \\
\rulename{mtruth}    & \ensuremath{{\mathcal V}_{\sigma}\llbracket\textbf{truth}\{\Varid{e}\}\rrbracket{}\pi}                  & = & \ensuremath{\textbf{truth}\{{\mathcal V}_{\sigma}\llbracket\Varid{e}\rrbracket{}\pi\}} \\
\rulename{mbind}   & \ensuremath{{\mathcal V}_{\sigma}\llbracket\exists{}\;\Varid{x}\rrbracket{}\pi}          & = & \ensuremath{\exists{}\;\Varid{x}} \\
\rulename{meq}     & \ensuremath{{\mathcal V}_{\sigma}\llbracket\Varid{e}_{\mathrm{1}}\mathrel{=}\Varid{e}_{\mathrm{2}}\rrbracket{}\pi}          & = & \ensuremath{{\mathcal V}_{\sigma}\llbracket\Varid{e}_{\mathrm{1}}\rrbracket{}\pi\mathrel{=}{\mathcal V}_{\sigma}\llbracket\Varid{e}_{\mathrm{2}}\rrbracket{}\pi} \\
\rulename{msemi}   & \ensuremath{{\mathcal V}_{\sigma}\llbracket\Varid{e}_{\mathrm{1}};\,\Varid{e}_{\mathrm{2}}\rrbracket{}\pi}        & = & \ensuremath{{\mathcal V}_{\sigma}\llbracket\Varid{e}_{\mathrm{1}}\rrbracket{}\pi;\,{\mathcal V}_{\sigma}\llbracket\Varid{e}_{\mathrm{2}}\rrbracket{}\pi} \\
\rulename{mchoice} & \ensuremath{{\mathcal V}_{\sigma}\llbracket\Varid{e}_{\mathrm{1}}\hspace{0.19em}\mathop{\rule[-0.1em]{0.15em}{0.75em}}\hspace{0.2em}\Varid{e}_{\mathrm{2}}\rrbracket{}\pi} & = & \ensuremath{{\mathcal V}_{\sigma}\llbracket\Varid{e}_{\mathrm{1}}\rrbracket{}\pi\hspace{0.19em}\mathop{\rule[-0.1em]{0.15em}{0.75em}}\hspace{0.2em}{\mathcal V}_{\sigma}\llbracket\Varid{e}_{\mathrm{2}}\rrbracket{}\pi} \\
\rulename{mfor}    & \ensuremath{{\mathcal V}_{\sigma}\llbracket\textbf{for}(\Varid{e}_{\mathrm{1}})\,\textbf{do}\,\Varid{e}_{\mathrm{2}}\rrbracket{}\pi}
   & = & \ensuremath{\textbf{iter}(\lambda\Varid{a}\;\Varid{r} .\,\textbf{cons}(\Varid{a}\langle\rangle,\Varid{r}\langle\rangle))\{{\mathcal V}_{\sigma}\llbracket\Varid{e}_{\mathrm{1}}\rrbracket{}\pi;\,\lambda{\bullet} .\,{\mathcal V}_{\sigma}\llbracket\Varid{e}_{\mathrm{2}}\rrbracket{}\pi\}\{\langle\rangle\}} \\
\rulename{mone}    & \ensuremath{{\mathcal V}_{\sigma}\llbracket\textbf{one}\{\Varid{e}\}\rrbracket{}\pi} & = & \ensuremath{\textbf{iter}(\lambda\Varid{a}\;{\bullet} .\,\Varid{a})\{{\mathcal V}_{\sigma}\llbracket\Varid{e}\rrbracket{}\pi\}\{\textbf{fail}\}} \\
\rulename{mall}    & \ensuremath{{\mathcal V}_{\sigma}\llbracket\textbf{all}\{\Varid{e}\}\rrbracket{}\pi} & = & \ensuremath{{\mathcal V}_{\sigma}\llbracket\textbf{for}(\exists\Varid{v};\,\Varid{v}\hspace{-0.14em}=\hspace{-0.14em}\Varid{e})\,\textbf{do}\,\Varid{v}\rrbracket{}\pi} \\
\multicolumn{2}{l}{\rulename{mif} \;\; \hfill \ensuremath{{\mathcal V}_{\sigma}\llbracket\;\textbf{if}\;\Varid{e}_{\mathrm{1}}\;\textbf{then}\;\Varid{e}_{\mathrm{2}}\;\textbf{else}\;\Varid{e}_{\mathrm{3}}\rrbracket{}\pi}}
& = & \multicolumn{2}{@{}l}{\ensuremath{\textbf{iter}(\lambda\Varid{a}\;{\bullet} .\,\Varid{a}\langle\rangle)\{{\mathcal V}_{\sigma}\llbracket\Varid{e}_{\mathrm{1}}\rrbracket{}\pi;\,\lambda{\bullet} .\,{\mathcal V}_{\sigma}\llbracket\Varid{e}_{\mathrm{2}}\rrbracket{}\pi\}\{{\mathcal V}_{\sigma}\llbracket\Varid{e}_{\mathrm{3}}\rrbracket{}\pi\}}}  \\
\rulename{mcheck+} & \ensuremath{{\mathcal V}_{+}\llbracket{\bf chk}\langle\omega\rangle\{\Varid{e}\}\rrbracket{}\pi} & = & \ensuremath{{\bf chk}\langle\omega\rangle\{{\mathcal V}_{+}\llbracket\Varid{e}\rrbracket{}\pi\}} \\
\rulename{mcheck-} & \ensuremath{{\mathcal V}_{-}\llbracket{\bf chk}\langle\omega\rangle\{\Varid{e}\}\rrbracket{}\pi} & = & \ensuremath{{\mathcal V}_{-}\llbracket\Varid{e}\rrbracket{}\pi} \\
\rulename{mcheckx} & \ensuremath{{\mathcal V}_{\textbf{x}}\llbracket{\bf chk}\langle\omega\rangle\{\Varid{e}\}\rrbracket{}\pi} & = & \ensuremath{{\bf chk}\langle\omega\rangle\{{\mathcal V}_{\textbf{x}}\llbracket\Varid{e}\rrbracket{}\pi\}} \\
\end{array}
\\

%--------- (e1 |> e2) and check<fx>{e}----------------
\begin{array}{lr@{\;\;}c@{\;\;}ll}
\multicolumn{4}{l}{\mbox{\it Type rules}} \\
\rulename{moftype+} & \ensuremath{{\mathcal V}_{+}\llbracket\Varid{e}_{\mathrm{1}}\vartriangleright\hspace{-1.3mm} \langle\omega\rangle\,\Varid{e}_{\mathrm{2}}\rrbracket{}\pi}
& = & \begin{array}[t]{@{}l}
  \ensuremath{\textbf{verify}} (;) \left\{ \begin{array}{@{}l}
       \ensuremath{\exists\Varid{r};\,\Varid{r}\hspace{-0.14em}=\hspace{-0.14em}{\bf chk}\langle\omega\rangle\{{\mathcal V}_{+}\llbracket\Varid{e}_{\mathrm{1}}\rrbracket{}\pi\};\,} \\
       \ensuremath{{\bf chk}\langle\textbf{succeeds}\rangle\{({\mathcal V}_{+}\llbracket\Varid{e}_{\mathrm{2}}\rrbracket{}\pi) [\!\Varid{r}]\})}
  \end{array} \right\} ; \\
  \ensuremath{{\mathcal V}_{-}\llbracket\Varid{e}_{\mathrm{1}}\vartriangleright\hspace{-1.3mm} \langle\omega\rangle\,\Varid{e}_{\mathrm{2}}\rrbracket{}\pi}
  \end{array} \\
\rulename{moftype-} & \ensuremath{{\mathcal V}_{-}\llbracket\Varid{e}_{\mathrm{1}}\vartriangleright\hspace{-1.3mm} \langle\omega\rangle\,\Varid{e}_{\mathrm{2}}\rrbracket{}\pi}
 & = & \ensuremath{\pi\;\!\blacktriangleright{}\!\;({\mathcal K}\langle{}\omega\rangle{};\,{\bf some}({\mathcal V}_{-}\llbracket\Varid{e}_{\mathrm{2}}\rrbracket{}\epsilon))} \\
\rulename{moftypex} & \ensuremath{{\mathcal V}_{\textbf{x}}\llbracket\Varid{e}_{\mathrm{1}}\vartriangleright\hspace{-1.3mm} \langle\omega\rangle\,\Varid{e}_{\mathrm{2}}\rrbracket{}\pi}
& = & \begin{array}[t]{@{}l}
    \ensuremath{\exists\Varid{r};\,\Varid{r}\hspace{-0.14em}=\hspace{-0.14em}{\bf chk}\langle\omega\rangle\{{\mathcal V}_{\textbf{x}}\llbracket\Varid{e}_{\mathrm{1}}\rrbracket{}\pi\};\,} \\
    \ensuremath{{\bf chk}\langle\textbf{succeeds}\rangle\{({\mathcal V}_{\textbf{x}}\llbracket\Varid{e}_{\mathrm{2}}\rrbracket{}\pi) [\!\Varid{r}]\})}
  \end{array}
\end{array}
\\
% -------------- Function rules ----------------
\begin{array}{lr@{\;\;}c@{\;\;}ll}
  \multicolumn{4}{l}{\mbox{\it Closed functions}} \\
 \rulename{mcfun+} &  \ensuremath{{\mathcal V}_{+}\llbracket\lambda_{\textbf{c}}\Varid{x}.(\Varid{e}_{\mathrm{1}})\{\Varid{e}_{\mathrm{2}}\}\rrbracket{}\pi}
 & = & \multicolumn{2}{@{}l}{\ensuremath{\textbf{verify}(\Varid{x};\,)\{{\mathcal V}_{-}\llbracket\Varid{e}_{\mathrm{1}}\rrbracket{}(\pi\mathbin{+}\Varid{x});\,{\mathcal V}_{+}\llbracket\Varid{e}_{\mathrm{2}}\rrbracket{}(\pi\mathbin{+}\Varid{x})\};\,}} \\
 & & &       \ensuremath{{\mathcal V}_{-}\llbracket\lambda_{\textbf{c}}\Varid{x}.(\Varid{e}_{\mathrm{1}})\{\Varid{e}_{\mathrm{2}}\}\rrbracket{}\pi} \\[1mm]
\rulename{mcfun-} &  \ensuremath{{\mathcal V}_{-}\llbracket\lambda_{\textbf{c}}\Varid{x}.(\Varid{e}_{\mathrm{1}})\{\Varid{e}_{\mathrm{2}}\}\rrbracket{}\pi}
& = & \ensuremath{\lambda\Varid{x} .\,{\mathcal V}_{+}\llbracket\Varid{e}_{\mathrm{1}}\rrbracket{}(\pi\mathbin{+}\Varid{x});\,{\mathcal V}_{-}\llbracket\Varid{e}_{\mathrm{2}}\rrbracket{}(\pi\mathbin{+}\Varid{x})} \\[1mm]
\rulename{mcfunx}& \ensuremath{{\mathcal V}_{\textbf{x}}\llbracket\lambda_{\textbf{c}}\Varid{x}.(\Varid{e}_{\mathrm{1}})\{\Varid{e}_{\mathrm{2}}\}\rrbracket{}\pi}
& = & \ensuremath{\lambda\Varid{x} .\,{\mathcal V}_{\textbf{x}}\llbracket\Varid{e}_{\mathrm{1}}\rrbracket{}(\pi\mathbin{+}\Varid{x});\,{\mathcal V}_{\textbf{x}}\llbracket\Varid{e}_{\mathrm{2}}\rrbracket{}(\pi\mathbin{+}\Varid{x})} \\[2mm]
\end{array} \\
% -------------- Auxiliary functions ----------------
\begin{array}{lr@{\;\;}c@{\;\;}ll}
\multicolumn{5}{l}{\mbox{{\bf Auxiliary functions:} $\ensuremath{{\mathcal K}\langle{}\omega\rangle{}\mathrel{=}\Varid{e}}$}} \\[0.02in]
\rulename{havoc-succ}     & \ensuremath{{\mathcal K}\langle{}\textbf{succeeds}\rangle{}\{\Varid{e}\}}  & = & \ensuremath{\Varid{e}} \\
\rulename{havoc-fail}     & \ensuremath{{\mathcal K}\langle{}\textbf{fails}\rangle{}\{\Varid{e}\}}  & = & \ensuremath{\textbf{fail}} \\
\rulename{havoc-dec}     & \ensuremath{{\mathcal K}\langle{}\textbf{decides}\rangle{}\{\Varid{e}\}}  & = & \ensuremath{{\bf some}(\lambda\Varid{x} .\,\Varid{x})\hspace{-0.14em}=\hspace{-0.14em}\langle{}\rangle;\,\Varid{e}} \\[1mm]
\end{array}

\end{array}$$
\caption{Desugaring Mini Verse into ICFP Verse} \label{fig:ds12}
\end{figure}


% ---------------------------------------------------------------------
%
%          The old, complicated DS12 Figure
%
%     Plus a second figure about open lambdas
%
% ----------------------------------------------------------------------

\begin{comment}
\begin{figure}
$$
\begin{array}{@{}l}
\mbox{\bf Syntax} \\
\multicolumn{1}{c}{\begin{array}{rcll}
  \ensuremath{\pi} & ::= &  \ensuremath{\mathsf{P}(\!\Varid{e})}  \vb \ensuremath{\mathsf{E}} \\
  \ensuremath{\sigma} & ::= & \ensuremath{\textbf{x}} \vb \ensuremath{+} \vb \ensuremath{-}  & \text{``Mode'': execution, verify, check} \\
\end{array}} \\

\begin{array}{lr@{\;\;}c@{\;\;}ll}
\multicolumn{5}{l}{\mbox{{\bf Auxiliary functions:} $\ensuremath{{\mathcal D}_{\sigma}\llbracket\Varid{t}\rrbracket} = \ensuremath{\Varid{e}}$, $\ensuremath{{\mathcal B}_{\sigma}\llbracket{}\Varid{t}\rrbracket{}\pi\,\Varid{j}} = \ensuremath{\Varid{e}}, \ensuremath{{\mathcal K}\langle{}\omega\rangle{}\mathrel{=}\Varid{e}}$}} \\[0.02in]
\rulename{dinit}   & \ensuremath{{\mathcal D}_{\sigma}\llbracket\Varid{t}\rrbracket} & = & \ensuremath{{\mathcal M}_{\sigma}\llbracket\Varid{t}\rrbracket{}\mathsf{E}} \\[1mm]
\rulename{dbody1}   & \ensuremath{{\mathcal B}_{\sigma}\llbracket{}\Varid{t}\rrbracket{}\mathsf{E}\,\Varid{j}} & = & \ensuremath{{\mathcal M}_{\sigma}\llbracket\Varid{t}\rrbracket{}\mathsf{E}}\\
\rulename{dbody2}   & \ensuremath{{\mathcal B}_{\sigma}\llbracket{}\Varid{t}\rrbracket{}\mathsf{P}(\!\Varid{f})\,\Varid{j}} & = & \ensuremath{{\mathcal M}_{\sigma}\llbracket\Varid{t}\rrbracket{}\mathsf{P}(\!\Varid{f} [\!\Varid{j}])}\\[1mm]
% \rulename{hide}     & |Hide t|  & = & |forall (r_1 r_2) (assume (r_1 == apply1N t r_2); r_1)| \\[1mm]
\rulename{funchk-e}     & \ensuremath{{\mathcal F}(\mathsf{E})} & = & \ensuremath{\langle{}\rangle} \\
\rulename{funchk-p}     & \ensuremath{{\mathcal F}(\mathsf{P}(\!\Varid{f}))} & = & \ensuremath{\textbf{isFun} [\!\Varid{f}]} \\[1mm]

\rulename{havoc-succ}     & \ensuremath{{\mathcal K}\langle{}\textbf{succeeds}\rangle{}\{\Varid{e}\}}  & = & \ensuremath{\Varid{e}} \\
\rulename{havoc-fail}     & \ensuremath{{\mathcal K}\langle{}\textbf{fails}\rangle{}\{\Varid{e}\}}  & = & \ensuremath{\textbf{fail}} \\
\rulename{havoc-dec}     & \ensuremath{{\mathcal K}\langle{}\textbf{decides}\rangle{}\{\Varid{e}\}}  & = & \ensuremath{{\bf some}(\lambda\Varid{x} .\,\Varid{x})\hspace{-0.14em}=\hspace{-0.14em}\langle{}\rangle;\,\Varid{e}} \\[1mm]
\end{array} \\

\begin{array}{lr@{\;}c@{\;}ll}
\multicolumn{5}{l}{\mbox{{\bf Desugaring expressions:} $\ensuremath{{\mathcal M}_{\sigma}\llbracket\Varid{t}\rrbracket{}\pi} = \ensuremath{\Varid{e}}$}} \\[0.02in]

\multicolumn{4}{l}{\mbox{\it Closed functions}} \\
 \rulename{mcfun+} &  \ensuremath{{\mathcal M}_{+}\llbracket\textbf{fun}_{\textbf{c}}(\Varid{t}_{\mathrm{1}})\{\Varid{t}_{\mathrm{2}}\}\rrbracket{}\pi}
 & = & \multicolumn{2}{@{}l}{\ensuremath{\textbf{verify}(\Varid{r};\,)\{\Varid{j}\mathrel{\coloneqq}{\mathcal M}_{-}\llbracket\Varid{t}_{\mathrm{1}}\rrbracket{}\mathsf{P}(\!\Varid{r});\,{\mathcal B}_{+}\llbracket{}\Varid{t}_{\mathrm{2}}\rrbracket{}\pi\,\Varid{j}\};\,}} \\
    &&& \ensuremath{{\mathcal M}_{-}\llbracket\textbf{fun}_{\textbf{c}}(\Varid{t}_{\mathrm{1}})\{\Varid{t}_{\mathrm{2}}\}\rrbracket{}\pi} \\[1mm]
\rulename{mcfun-} &  \ensuremath{{\mathcal M}_{-}\llbracket\textbf{fun}_{\textbf{c}}(\Varid{t}_{\mathrm{1}})\{\Varid{t}_{\mathrm{2}}\}\rrbracket{}\pi}
& = & \ensuremath{{\mathcal F}(\pi);\,\lambda\Varid{i} .\,\Varid{j}\mathrel{\coloneqq}{\mathcal M}_{+}\llbracket\Varid{t}_{\mathrm{1}}\rrbracket{}\mathsf{P}(\!\Varid{i});\,{\mathcal B}_{-}\llbracket{}\Varid{t}_{\mathrm{2}}\rrbracket{}\pi\,\Varid{j}} \\[1mm]
\rulename{mcfunx}& \ensuremath{{\mathcal M}_{\textbf{x}}\llbracket\textbf{fun}_{\textbf{c}}(\Varid{t}_{\mathrm{1}})\{\Varid{t}_{\mathrm{2}}\}\rrbracket{}\pi}
& = & \ensuremath{{\mathcal F}(\pi);\,\lambda\Varid{i} .\,\Varid{j}\mathrel{\coloneqq}{\mathcal M}_{\textbf{x}}\llbracket\Varid{t}_{\mathrm{1}}\rrbracket{}\mathsf{P}(\!\Varid{i});\,{\mathcal B}_{\textbf{x}}\llbracket{}\Varid{t}_{\mathrm{2}}\rrbracket{}\pi\,\Varid{j}} \\[2mm]
\end{array}
\\
\begin{array}{lr@{\;\;}c@{\;\;}ll}

%--------- (t1 |><fx> t2) ----------------
\multicolumn{4}{l}{\mbox{\it Type rules}} \\
\rulename{moftype-} & \ensuremath{{\mathcal M}_{-}\llbracket\Varid{t}_{\mathrm{1}}\vartriangleright\hspace{-1.3mm} \langle\omega\rangle\,\Varid{t}_{\mathrm{2}}\rrbracket{}\pi}
  & = & \ensuremath{\Varid{z}\mathrel{\coloneqq}{\mathcal D}_{-}\llbracket\Varid{t}_{\mathrm{2}}\rrbracket;\,\mathsf{fvs}(\Varid{t}_{\mathrm{1}})\;\!\blacktriangleright{}\!\;({\mathcal K}\langle{}\omega\rangle{}\{{\bf some}(\Varid{z})\})} \\
\rulename{moftype+x} & \ensuremath{{\mathcal M}_{\sigma}\llbracket\Varid{t}_{\mathrm{1}}\vartriangleright\hspace{-1.3mm} \langle\omega\rangle\,\Varid{t}_{\mathrm{2}}\rrbracket{}\pi} & = &
  \begin{array}[t]{@{}l}
  \ensuremath{\Varid{y}\mathrel{\coloneqq}{\bf chk}\langle\omega\rangle\{{\mathcal M}_{\textbf{x}}\llbracket\Varid{t}_{\mathrm{1}}\rrbracket{}\pi\};\,} \\
  \ensuremath{{\bf chk}\langle\textbf{succeeds}\rangle\{\Varid{z}\mathrel{\coloneqq}{\mathcal D}_{\textbf{x}}\llbracket\Varid{t}_{\mathrm{2}}\rrbracket;\,\Varid{z} [\!\Varid{y}]\}}
  \end{array}   & \ensuremath{\sigma}\in\{\ensuremath{+},\ensuremath{\textbf{x}}\} \\

%--------- (:<fx> t) ----------------
\rulename{mtype-}  & \ensuremath{{\mathcal M}_{-}\llbracket\mathbin{:}\Varid{t}\hspace{-0.2mm} \langle\omega\rangle\rrbracket{}\mathsf{P}(\!\Varid{e})} & = & \ensuremath{\Varid{z}\mathrel{\coloneqq}{\mathcal D}_{-}\llbracket\Varid{t}\rrbracket;\,\mathsf{fvs}(\Varid{e})\;\!\blacktriangleright{}\!\;({\mathcal K}\langle{}\omega\rangle{}\{{\bf some}(\Varid{z})\})} \\
\rulename{mtype+x1} & \ensuremath{{\mathcal M}_{\sigma}\llbracket\mathbin{:}\Varid{t}\hspace{-0.2mm} \langle \rangle\rrbracket{}\mathsf{P}(\!\Varid{e})} & = & \ensuremath{\Varid{y}\mathrel{\coloneqq}\Varid{e};\,\Varid{z}\mathrel{\coloneqq}{\mathcal D}_{\sigma}\llbracket\Varid{t}\rrbracket;\,\Varid{z} [\!\Varid{y}]}  & \ensuremath{\sigma}\in\{\ensuremath{+},\ensuremath{\textbf{x}}\}\\
\rulename{mtype+x2} & \ensuremath{{\mathcal M}_{\sigma}\llbracket\mathbin{:}\Varid{t}\hspace{-0.2mm} \langle\omega\rangle\rrbracket{}\mathsf{P}(\!\Varid{e})} & = &
  \begin{array}[t]{@{}l}
  \ensuremath{\Varid{y}\mathrel{\coloneqq}{\bf chk}\langle\omega\rangle\{\Varid{e}\};\,} \\
  \ensuremath{{\bf chk}\langle\textbf{succeeds}\rangle\{\Varid{z}\mathrel{\coloneqq}{\mathcal D}_{\textbf{x}}\llbracket\Varid{t}\rrbracket;\,\Varid{z} [\!\Varid{y}]\}}
  \end{array}   & \ensuremath{\sigma}\in\{\ensuremath{+},\ensuremath{\textbf{x}}\} \\
\rulename{mtypee}  & \ensuremath{{\mathcal M}_{\sigma}\llbracket\mathbin{:}\Varid{t}\hspace{-0.2mm} \langle\omega\rangle\rrbracket{}\mathsf{E}} & = & \ensuremath{\exists\Varid{x}.\,\Varid{z}\mathrel{\coloneqq}{\mathcal D}_{\sigma}\llbracket\Varid{t}\rrbracket;\,\Varid{z} [\!\Varid{x}]} \\

%--------- check<fx>{t} ----------------
\rulename{mcheck-}  & \ensuremath{{\mathcal M}_{-}\llbracket{\bf chk}\langle\omega\rangle\{\Varid{t}\}\rrbracket{}\pi} & = & \ensuremath{{\mathcal M}_{-}\llbracket\Varid{t}\rrbracket{}\pi} \\
\rulename{mcheck+x} & \ensuremath{{\mathcal M}_{\sigma}\llbracket{\bf chk}\langle\omega\rangle\{\Varid{t}\}\rrbracket{}\pi} & = & \ensuremath{{\bf chk}\langle\omega\rangle\{{\mathcal M}_{\sigma}\llbracket\Varid{t}\rrbracket{}\pi\}} & \ensuremath{\sigma}\in\{\ensuremath{+},\ensuremath{\textbf{x}}\} \\

\rulename{msquige} & \ensuremath{{\mathcal M}_{\sigma}\llbracket\Varid{x}\!\rightarrow\!\Varid{y}\mathrel{\coloneqq}\Varid{t}\rrbracket{}\mathsf{E}}
     & = & \ensuremath{\exists\Varid{i}.\,{\mathcal M}_{\sigma}\llbracket\Varid{x}\!\rightarrow\!\Varid{y}\mathrel{\coloneqq}\Varid{t}\rrbracket{}\mathsf{P}(\!\Varid{i})} \\
\rulename{msquigp} & \ensuremath{{\mathcal M}_{\sigma}\llbracket\Varid{x}\!\rightarrow\!\Varid{y}\mathrel{\coloneqq}\Varid{t}\rrbracket{}\mathsf{P}(\!\Varid{i})}
     & = & \ensuremath{\Varid{x}\mathrel{\coloneqq}\Varid{i};\,\Varid{y}\mathrel{\coloneqq}{\mathcal M}_{\sigma}\llbracket\Varid{t}\rrbracket{}\mathsf{P}(\!\Varid{i})} \\
\rulename{marraye} & \ensuremath{{\mathcal M}_{\sigma}\llbracket\langle\Varid{t}_{\mathrm{1}},\mydots,\Varid{t}_{\Varid{n}}\rangle\rrbracket{}\mathsf{E}} & = &
    \ensuremath{\langle{\mathcal M}_{\sigma}\llbracket\Varid{t}_{\mathrm{1}}\rrbracket{}\mathsf{E},\mydots,{\mathcal M}_{\sigma}\llbracket\Varid{t}_{\Varid{n}}\rrbracket{}\mathsf{E}\rangle} \\
\rulename{marrayp} & \ensuremath{{\mathcal M}_{\sigma}\llbracket\langle\Varid{t}_{\mathrm{1}},\mydots,\Varid{t}_{\Varid{n}}\rangle\rrbracket{}\mathsf{P}(\!\Varid{i})} & = &
    \multicolumn{2}{@{}l}{\ensuremath{\exists\Varid{i}_{\mathrm{1}}\;\mydots\;\Varid{i}_{\Varid{n}}.\,\Varid{i}\hspace{-0.14em}=\hspace{-0.14em}\langle\Varid{i}_{\mathrm{1}},\mydots,\Varid{i}_{\Varid{n}}\rangle;\,\langle{\mathcal M}_{\sigma}\llbracket\Varid{t}_{\mathrm{1}}\rrbracket{}\mathsf{P}(\!\Varid{i}_{\mathrm{1}}),\mydots,{\mathcal M}_{\sigma}\llbracket\Varid{t}_{\Varid{n}}\rrbracket{}\mathsf{P}(\!\Varid{i}_{\Varid{n}})\rangle}}
    \\[2mm]

\end{array}
\\
\begin{array}{lr@{\;\;}c@{\;\;}ll}
\multicolumn{4}{l}{\mbox{\it Propagation rules}} \\
\rulename{mbind}   & \ensuremath{{\mathcal M}_{\sigma}\llbracket\Varid{x}\mathrel{\coloneqq}\Varid{t}\rrbracket{}\pi}         & = & \ensuremath{\Varid{x}\mathrel{\coloneqq}{\mathcal M}_{\sigma}\llbracket\Varid{t}\rrbracket{}\pi} \\
\rulename{meq}     & \ensuremath{{\mathcal M}_{\sigma}\llbracket\Varid{t}_{\mathrm{1}}\mathrel{=}\Varid{t}_{\mathrm{2}}\rrbracket{}\pi}      & = & \ensuremath{{\mathcal M}_{\sigma}\llbracket\Varid{t}_{\mathrm{1}}\rrbracket{}\pi\mathrel{=}{\mathcal M}_{\sigma}\llbracket\Varid{t}_{\mathrm{2}}\rrbracket{}\pi} \\
\rulename{msemi}   & \ensuremath{{\mathcal M}_{\sigma}\llbracket\Varid{t}_{\mathrm{1}};\,\Varid{t}_{\mathrm{2}}\rrbracket{}\pi}    & = & \ensuremath{{\mathcal D}_{\sigma}\llbracket\Varid{t}_{\mathrm{1}}\rrbracket;\,{\mathcal M}_{\sigma}\llbracket\Varid{t}_{\mathrm{2}}\rrbracket{}\pi} \\
\rulename{mchoice} & \ensuremath{{\mathcal M}_{\sigma}\llbracket\Varid{t}_{\mathrm{1}}\hspace{0.19em}\mathop{\rule[-0.1em]{0.15em}{0.75em}}\hspace{0.2em}\Varid{t}_{\mathrm{2}}\rrbracket{}\pi} & = & \ensuremath{{\mathcal M}_{\sigma}\llbracket\Varid{t}_{\mathrm{1}}\rrbracket{}\pi\hspace{0.19em}\mathop{\rule[-0.1em]{0.15em}{0.75em}}\hspace{0.2em}{\mathcal M}_{\sigma}\llbracket\Varid{t}_{\mathrm{2}}\rrbracket{}\pi} \\
\rulename{mif} & \ensuremath{{\mathcal M}_{\sigma}\llbracket\;\textbf{if}\;\Varid{t}_{\mathrm{1}}\;\textbf{then}\;\Varid{t}_{\mathrm{2}}\;\textbf{else}\;\Varid{t}_{\mathrm{3}}\rrbracket{}\pi}
  & = & \multicolumn{2}{@{}l}{\ensuremath{\textbf{if}\;{\mathcal D}_{\sigma}\llbracket\Varid{t}_{\mathrm{1}}\rrbracket\;\textbf{then}\;{\mathcal M}_{\sigma}\llbracket\Varid{t}_{\mathrm{2}}\rrbracket{}\pi\;\textbf{else}\;{\mathcal M}_{\sigma}\llbracket\Varid{t}_{\mathrm{3}}\rrbracket{}\pi}} \\[2mm]

\multicolumn{4}{l}{\mbox{\it Equality rules}} \\
\rulename{mconst}  & \ensuremath{{\mathcal M}_{\sigma}\llbracket\Varid{k}\rrbracket{}\mathsf{E}}                  & = & \ensuremath{\Varid{k}} \\
nnn\rulename{mvar}    & \ensuremath{{\mathcal M}_{\sigma}\llbracket\Varid{y}\rrbracket{}\mathsf{E}}                  & = & \ensuremath{\Varid{y}} \\
\rulename{mapp}    & \ensuremath{{\mathcal M}_{\sigma}\llbracket\Varid{t}_{\mathrm{1}} [\!\Varid{t}_{\mathrm{2}}]\rrbracket{}\mathsf{E}}  & = & \ensuremath{\Varid{f}\mathrel{\coloneqq}{\mathcal D}_{\sigma}\llbracket\Varid{t}_{\mathrm{1}}\rrbracket;\,\Varid{a}\mathrel{\coloneqq}{\mathcal D}_{\sigma}\llbracket\Varid{t}_{\mathrm{2}}\rrbracket;\,\Varid{f} [\!\Varid{a}]} \\
\rulename{meq}     & \ensuremath{{\mathcal M}_{\sigma}\llbracket\Varid{t}\rrbracket{}\mathsf{P}(\!\Varid{e})} & = & \ensuremath{{\mathcal M}_{\sigma}\llbracket\Varid{t}\rrbracket{}\mathsf{E}\mathrel{=}\Varid{e}}
\end{array}
\end{array}$$
\caption{{\bf DS12}: Small Source to Big Core } \label{fig:ds12}
\end{figure}

\begin{figure}[t]
  $$
\begin{array}{@{}l}
\begin{array}{@{}lr@{\;\;}c@{\;\;}ll}
\multicolumn{4}{@{}l}{\mbox{\it Desugaring rules for open functions}} \\
\rulename{mofun+}  & \ensuremath{{\mathcal M}\llbracket\textbf{fun}_{\textbf{o}}(\Varid{t}_{\mathrm{1}})\langle\omega\rangle\{\Varid{t}_{\mathrm{2}}\}\rrbracket{}\Varid{f}^{\pi}_{+}}
& = & \multicolumn{2}{@{}l}{\ensuremath{\textbf{verify}\{\forall\Varid{i}.\,\Varid{j}\mathrel{\coloneqq}\Me{ \Varid{i}^{\circ}_{-}}{ \Varid{t}_{\mathrm{1}}};\,{\bf chk}\langle\omega\rangle\{\Varid{z}\mathrel{\coloneqq}\Conid{DAe}\;(\Varid{f}^{\pi})\;\Varid{j};\,\Me{ \Varid{z}^{\pi}}{ \Varid{t}_{\mathrm{2}}}\}\};\,}} \\
& &   & \ensuremath{{\mathcal M}\llbracket\textbf{fun}_{\textbf{o}}(\Varid{t}_{\mathrm{1}})\langle\omega\rangle\{\Varid{t}_{\mathrm{2}}\}\rrbracket{}\Varid{f}^{\pi}_{-}} \\[1mm]

\rulename{mofun-}  & \ensuremath{{\mathcal M}\llbracket\textbf{fun}_{\textbf{o}}(\Varid{t}_{\mathrm{1}})\langle\omega\rangle\{\Varid{t}_{\mathrm{2}}\}\rrbracket{}\Varid{f}^{\pi}_{-}}
 & = & \ensuremath{\exists\Varid{h}.\,\;\textbf{olam}} \left( \begin{array}{@{}l@{}}
    \ensuremath{\Varid{h},} \\
    \ensuremath{\lambda\Varid{i} .\,\Varid{j}\mathrel{\coloneqq}\Me{ \Varid{i}^{\bullet}_{\textbf{x}}}{ \Varid{t}_{\mathrm{1}}};\,\langle\Varid{j},\Varid{x}_{\mathrm{1}},..,\Varid{x}_{\Varid{n}}\rangle,} \\
    \ensuremath{\lambda\Varid{x}.\,} \begin{array}[t]{@{}l@{}}
      \ensuremath{(\langle\Varid{j},\Varid{x}_{\mathrm{1}},..,\Varid{x}_{\Varid{n}}\rangle\mathrel{\coloneqq}\Varid{x};\,{\mathcal K}\langle{}\omega\rangle{};\,} \\
      \ensuremath{\Varid{z}\mathrel{\coloneqq}\Conid{DAe}\;(\Varid{f}^{\pi})\;\Varid{j};\,\Me{ \Varid{z}^{\pi}_{-}}{ \Varid{t}_{\mathrm{2}}})}
      \end{array}
   \end{array} \right) \\
    && \multicolumn{2}{l}{\text{where \ensuremath{\Varid{x}_{\mathrm{1}},..,\Varid{x}_{\Varid{n}}} are the variables bound by \ensuremath{\Varid{t}_{\mathrm{1}}}}} \\

\rulename{mofunx}  & \ensuremath{{\mathcal M}\llbracket\textbf{fun}_{\textbf{o}}(\Varid{t}_{\mathrm{1}})\langle\omega\rangle\{\Varid{t}_{\mathrm{2}}\}\rrbracket{}\Varid{f}^{\pi}_{\textbf{x}}}
  & = & \ensuremath{\exists\Varid{h}.\,\;\textbf{olam}} \left( \begin{array}{@{}l@{}}
    \ensuremath{\Varid{h},} \\
    \ensuremath{\lambda\Varid{i} .\,\Varid{j}\mathrel{\coloneqq}\Me{ \Varid{i}^{\bullet}_{\textbf{x}}}{ \Varid{t}_{\mathrm{1}}};\,\langle\Varid{j},\Varid{x}_{\mathrm{1}},..,\Varid{x}_{\Varid{n}}\rangle,} \\
    \ensuremath{\lambda\Varid{x} .\,\langle\Varid{j},\Varid{x}_{\mathrm{1}},..,\Varid{x}_{\Varid{n}}\rangle\mathrel{\coloneqq}\Varid{x};\,\Varid{z}\mathrel{\coloneqq}\Conid{DAe}\;(\Varid{f}^{\pi})\;\Varid{j};\,\Me{ \Varid{z}^{\pi}_{\textbf{x}}}{ \Varid{t}_{\mathrm{2}}}}
   \end{array} \right) \\
    && \multicolumn{2}{l}{\text{where \ensuremath{\Varid{x}_{\mathrm{1}},..,\Varid{x}_{\Varid{n}}} are the variables bound by \ensuremath{\Varid{t}_{\mathrm{1}}}}}
\end{array}
\end{array}
$$
\caption{Extra stuff for open functions} \label{fig:open}
\end{figure}
\end{comment}

